% Contenu partagé — preuves, blocs étiquetés, mises en valeur.

\chapter{Preuves et blocs annexes}%
\label{chap:blocks}
\minitoc%

%---------------------------------------------------------------------------
\section{Preuves}%
\label{sec:proofs}

\begin{proposition}{Une proposition à démontrer}{ademontrer}
    Pour tout $x \in \R$, $x^2 \ge 0$.
\end{proposition}

\begin{proof}
    Le marqueur d'ouverture et le carré de fin viennent du thème, comme la
    couleur du texte. \newstep{} Un astérisque sépare deux étapes.
\end{proof}

\begin{proof}[Démonstration de la proposition~\ref{prop:ademontrer}]
    L'argument facultatif remplace le marqueur d'ouverture.
\end{proof}

Une preuve peut être étalée sur plusieurs diapositives. Sur papier, les trois
morceaux se suivent~; en diapositives, seul le dernier pose le carré final.

\begin{proofbegin}
    Premier morceau.
\end{proofbegin}
\begin{proofmiddle}
    Morceau intermédiaire.
\end{proofmiddle}
\begin{proofend}
    Dernier morceau.
\end{proofend}

%---------------------------------------------------------------------------
\section{Blocs étiquetés}%
\label{sec:tagged}

Trois familles numérotées indépendamment, marquées par un filet latéral ou un
simple retrait. Elles partagent la même fabrique interne~: en ajouter une
quatrième tient en une ligne.

\begin{assumption}
    Une hypothèse. Elles sont étiquetées H1, H2, \ldots
\end{assumption}

\begin{assumption}
    Une deuxième hypothèse.
\end{assumption}

\begin{assumption*}
    Une hypothèse non étiquetée.
\end{assumption*}

\begin{openquestion}
    Un questionnement à se poser. Étiquetage Q1, Q2, \ldots
\end{openquestion}

\begin{difficulty}
    Une difficulté attendue. Étiquetage D1, D2, \ldots
\end{difficulty}

\begin{web}[Documentation en ligne du cours]
    Renvoi vers une ressource extérieure, précédé de son pictogramme.
\end{web}

%---------------------------------------------------------------------------
\section{Mise en valeur}%
\label{sec:emphasis}

Le \keyword{terme important} se compose avec \verb|\keyword|. Les quatre
couleurs du thème sont accessibles~: \emphA{première}, \emphB{deuxième},
\emphC{troisième}, \emphD{quatrième}, plus \emphLink{la couleur des liens}.

Pictogrammes~: \cmark{} et \xmark{}, plus la main \handwrite{} qui précédait les
exercices en v0.

\supplement{Un complément, en vert.}

\reminder{Un rappel, en rouge.}

Une équation mise en valeur par un surlignage, dont la couleur vient du thème~:
\[
    \tcbhighmath{\dot{x}(t) = A\, x(t) + b(t)}
\]

Numérotation indépendante des problèmes~:
\begin{equation}
    \min_{u} \int_0^{t_f} \! \ell(x(t), u(t)) \dif t
    \tagProblem
\end{equation}

%---------------------------------------------------------------------------
\section{Notes de travail}%
\label{sec:draft}

Visibles seulement avec l'option \texttt{draft}, invisibles sinon. Elles
remplacent le \verb|\ifVIEWALL| de la v0, qui n'était jamais défini et
provoquait une erreur à l'usage.

\notework{Une note personnelle.}

\worktodo{Quelque chose à finir.}

\begin{worknotes}
    Un bloc entier de notes de travail, exclu de la compilation finale.
\end{worknotes}
